% =====================================================================================
% CERTA -- Supplementary Information
% Compile AFTER main.tex (main.tex declares \myexternaldocument{SI}; SI declares the reverse)
% so that Supplementary Figure/Table numbers resolve as cross-references in the main text.
% Requires: figS1.pdf..figS6.pdf (build_si_figures.py) and si_tables.tex (si_tables.py).
% Architecture: 12 sections, 6 supplementary figures, 16 supplementary tables.
% =====================================================================================
\documentclass[11pt]{article}
\usepackage[margin=1in]{geometry}
\usepackage{amsmath,amssymb}
\usepackage{graphicx}
\usepackage{booktabs}
\usepackage{makecell}
\usepackage{array}
\usepackage{caption}
\usepackage{microtype}
\usepackage{float}
\usepackage{amsthm}
\usepackage{xr-hyper}
\usepackage{hyperref}
\externaldocument{main}

% Supplementary numbering: S1, S2, ...
\renewcommand{\thefigure}{S\arabic{figure}}
\renewcommand{\thetable}{S\arabic{table}}
\renewcommand{\theequation}{S\arabic{equation}}
\renewcommand{\thesection}{S\arabic{section}}
\makeatletter\renewcommand{\l@section}{\@dottedtocline{1}{0pt}{2.7em}}\makeatother
\captionsetup{labelfont=bf,font=small}
\newfloat{algorithm}{tbp}{loa}
\floatname{algorithm}{Algorithm}
\renewcommand{\thealgorithm}{S\arabic{algorithm}}
\theoremstyle{plain}
\newtheorem{proposition}{Proposition}
\renewcommand{\theproposition}{S\arabic{proposition}}
\newcommand{\YES}{\mathrm{YES}}
\newcommand{\NO}{\mathrm{NO}}
\newcommand{\Prob}{\mathbb{P}}
\newcommand{\Risk}{R}

\newcommand{\tableSdecdist}{%
Qwen2.5-72B & 41.2 & 58.8 & 42.8 & 57.2 & 45.0/37.4 & 46.2/39.4 \\
HuatuoGPT-o1-72B & 40.5 & 59.5 & 41.9 & 58.1 & 42.8/38.2 & 43.6/40.2 \\
Med42-70B & 62.1 & 37.9 & 62.5 & 37.5 & 83.4/40.9 & 84.2/40.8 \\
OpenBioLLM-70B & 91.1 & 8.9 & 93.2 & 6.8 & 99.4/82.8 & 100.0/86.5 \\
MMed-Llama-3.1-70B & 86.8 & 13.2 & 89.0 & 11.0 & 100.0/73.6 & 100.0/77.9 \\
Llama-3.3-70B & 57.5 & 42.5 & 67.2 & 32.8 & 83.4/31.7 & 92.2/42.1 \\
Llama-3-70B (base) & 60.3 & 39.7 & 64.3 & 35.7 & 83.3/37.3 & 89.9/38.7 \\
Llama-3-8B & 66.0 & 34.0 & 70.4 & 29.6 & 98.0/34.0 & 98.6/42.2 \\
MMed-Llama-3-8B & 74.0 & 26.0 & 67.7 & 32.3 & 99.4/48.5 & 99.3/36.1 \\
Mistral-7B & 51.8 & 48.2 & 47.8 & 52.2 & 72.3/31.4 & 67.5/28.1 \\
BioMistral-7B & 94.3 & 5.7 & 99.5 & 0.5 & 100.0/88.6 & 100.0/99.0 \\
}
\newcommand{\tableSperconfig}{%
Qwen2.5-72B & 100.0\,[100.0,100.0] & 100.0\,[100.0,100.0] & 93.6\,[92.1,94.9] & 100.0\,[100.0,100.0] & 100.0\,[100.0,100.0] \\
HuatuoGPT-o1-72B & 100.0\,[100.0,100.0] & 100.0\,[100.0,100.0] & 91.0\,[89.5,92.6] & 100.0\,[100.0,100.0] & 100.0\,[100.0,100.0] \\
Med42-70B & 99.8\,[99.6,100.0] & 100.0\,[100.0,100.0] & 56.3\,[53.7,59.2] & 0.0\,[0.0,0.0] & 97.5\,[96.6,98.4] \\
OpenBioLLM-70B & 100.0\,[100.0,100.0] & 100.0\,[100.0,100.0] & 84.7\,[82.6,86.6] & 86.9\,[85.0,88.9] & 98.1\,[97.3,98.8] \\
MMed-Llama-3.1-70B & 100.0\,[100.0,100.0] & 93.6\,[92.2,95.0] & 72.1\,[69.6,74.4] & 74.6\,[72.1,77.1] & 96.9\,[95.9,97.8] \\
Llama-3.3-70B & 100.0\,[100.0,100.0] & 100.0\,[100.0,100.0] & 51.9\,[49.2,54.7] & 38.1\,[35.4,40.8] & 82.2\,[80.0,84.2] \\
Llama-3-70B (base) & 100.0\,[100.0,100.0] & 98.2\,[97.4,98.9] & 53.8\,[51.2,56.6] & 0.0\,[0.0,0.0] & 93.0\,[91.6,94.6] \\
Llama-3-8B & 100.0\,[100.0,100.0] & 100.0\,[100.0,100.0] & 15.0\,[13.0,17.1] & 100.0\,[100.0,100.0] & 92.8\,[91.4,94.2] \\
MMed-Llama-3-8B & 100.0\,[100.0,100.0] & 98.2\,[97.4,99.0] & 0.0\,[0.0,0.0] & 89.0\,[87.4,90.8] & 84.2\,[82.1,86.2] \\
Mistral-7B & 100.0\,[100.0,100.0] & 100.0\,[100.0,100.0] & 66.6\,[63.9,69.1] & 100.0\,[100.0,100.0] & 83.8\,[81.8,85.9] \\
BioMistral-7B & 100.0\,[100.0,100.0] & 80.3\,[78.2,82.6] & 56.1\,[53.4,58.9] & 100.0\,[100.0,100.0] & 95.8\,[94.6,96.8] \\
}
\newcommand{\tableSaudit}{%
Qwen2.5-72B & 1250 & 26 & 0.00 & 3.5 & 100.0 [100.0, 100.0] \\
HuatuoGPT-o1-72B & 1250 & 33 & 0.00 & 4.2 & 100.0 [100.0, 100.0] \\
Med42-70B & 1224 & 35 & 3.06 & 4.5 & 97.5 [96.6, 98.4] \\
OpenBioLLM-70B & 1224 & 37 & 0.31 & 4.7 & 98.1 [97.3, 98.8] \\
MMed-Llama-3.1-70B & 1223 & 29 & 0.12 & 3.9 & 96.9 [95.9, 97.8] \\
Llama-3.3-70B & 1048 & 32 & 8.25 & 4.9 & 82.2 [80.0, 84.3] \\
Llama-3-70B (base) & 1151 & 34 & 3.50 & 4.7 & 93.0 [91.5, 94.3] \\
Llama-3-8B & 1149 & 30 & 1.12 & 4.3 & 92.8 [91.4, 94.2] \\
MMed-Llama-3-8B & 1046 & 30 & 0.61 & 4.7 & 84.2 [82.4, 86.2] \\
Mistral-7B & 1074 & 25 & 0.12 & 4.0 & 83.8 [81.9, 85.8] \\
BioMistral-7B & 1206 & 32 & 0.38 & 4.3 & 95.8 [94.6, 96.9] \\
}
\newcommand{\tableSmond}{%
Qwen2.5-72B & 517 & 11 & 0.34 & 97.4 & 4.9 & 722 & 6 & 0.12 & 100.0 & 2.5 \\
HuatuoGPT-o1-72B & 480 & 10 & 0.54 & 92.7 & 5.0 & 738 & 8 & 0.12 & 100.0 & 2.8 \\
Med42-70B & 762 & 18 & 3.71 & 97.7 & 4.6 & 433 & 4 & 10.78 & 88.2 & 3.4 \\
OpenBioLLM-70B & 1140 & 29 & 0.38 & 98.7 & 4.3 & 0 & 0 & $\infty$ & \emph{esc} & -- \\
MMed-Llama-3.1-70B & 1078 & 21 & 0.12 & 96.5 & 3.6 & 0 & 0 & $\infty$ & \emph{esc} & -- \\
Llama-3.3-70B & 612 & 13 & 10.62 & 72.3 & 4.6 & 420 & 0 & 0.25 & 100.0 & 1.6 \\
Llama-3-70B (base) & 677 & 17 & 5.75 & 85.9 & 5.0 & 426 & 7 & 3.25 & 92.3 & 4.5 \\
Llama-3-8B & 775 & 20 & 1.62 & 89.6 & 4.9 & 370 & 3 & 0.62 & 94.8 & 3.5 \\
MMed-Llama-3-8B & 832 & 9 & 0.00 & 100.0 & 2.7 & 162 & 0 & 1.42 & 35.1 & 4.2 \\
Mistral-7B & 528 & 5 & 0.06 & 89.4 & 3.1 & 554 & 6 & 0.19 & 80.6 & 3.2 \\
BioMistral-7B & 1205 & 32 & 0.38 & 96.3 & 4.4 & 0 & 0 & $\infty$ & \emph{esc} & -- \\
}
\newcommand{\tableSclassres}{%
Qwen2.5-72B & 528 & 722 & 1.30 & 0.95 & 1.0 & certified \\
HuatuoGPT-o1-72B & 512 & 738 & 1.34 & 0.93 & 1.1 & certified \\
Med42-70B & 778 & 472 & 0.88 & 1.45 & 1.7 & certified \\
OpenBioLLM-70B & 1159 & 91 & 0.59 & 7.31 & 13.7 & \emph{resolution-limited} \\
MMed-Llama-3.1-70B & 1105 & 145 & 0.62 & 4.65 & 8.4 & \emph{risk-limited} \\
Llama-3.3-70B & 830 & 420 & 0.83 & 1.63 & 1.6 & certified \\
Llama-3-70B (base) & 787 & 463 & 0.87 & 1.48 & 1.8 & certified \\
Llama-3-8B & 857 & 393 & 0.80 & 1.74 & 2.0 & certified \\
MMed-Llama-3-8B & 832 & 418 & 0.83 & 1.64 & 4.2 & certified \\
Mistral-7B & 590 & 660 & 1.16 & 1.04 & 1.4 & certified \\
BioMistral-7B & 1246 & 4 & 0.55 & 82.22 & 82.2 & \emph{resolution-limited} \\
}
\newcommand{\tableSgrid}{%
25 & 25.0 & 29.4 & 37.5 & 25.0 & 25.0 & 37.5 \\
50 & 25.0 & 28.2 & 37.5 & 25.0 & 25.0 & 37.5 \\
100 & 25.0 & 28.2 & 38.2 & 25.0 & 25.0 & 38.2 \\
200 & 25.0 & 28.2 & 38.2 & 25.0 & 25.0 & 38.2 \\
}
\newcommand{\tableSsplit}{%
Qwen2.5-72B & 25.0 & 25.0--25.0 & 25.0 & 25.0 \\
HuatuoGPT-o1-72B & 25.0 & 25.0--25.0 & 25.0 & 25.8 \\
Med42-70B & 27.8 & 26.8--28.3 & 25.0 & 29.3 \\
OpenBioLLM-70B & 27.2 & 26.6--27.6 & 25.7 & 27.9 \\
MMed-Llama-3.1-70B & 26.7 & 25.0--26.9 & 25.0 & 27.3 \\
Llama-3.3-70B & 37.5 & 36.8--37.5 & 35.8 & 39.1 \\
Llama-3-70B (base) & 33.2 & 31.9--34.4 & 30.0 & 35.7 \\
Llama-3-8B & 29.9 & 29.2--30.3 & 27.6 & 31.5 \\
MMed-Llama-3-8B & 36.8 & 35.9--37.5 & 34.2 & 39.0 \\
Mistral-7B & 36.7 & 36.2--37.0 & 34.9 & 40.3 \\
BioMistral-7B & 27.9 & 26.9--28.0 & 26.6 & 28.7 \\
}
\newcommand{\tableSconfidence}{%
Qwen2.5-72B & 25.0 & 25.0 & 25.0 & 25.8 & 26.2 & 26.2 \\
HuatuoGPT-o1-72B & 25.0 & 25.0 & 25.0 & 27.3 & 27.6 & 28.2 \\
Med42-70B & 26.9 & 26.9 & 26.9 & 29.4 & 29.4 & 30.6 \\
OpenBioLLM-70B & 26.4 & 26.4 & 27.6 & 30.6 & 30.6 & 30.6 \\
MMed-Llama-3.1-70B & 27.3 & 27.3 & 27.3 & 35.2 & 35.2 & 35.2 \\
Llama-3.3-70B & 37.5 & 38.3 & 38.3 & 39.1 & 39.1 & 40.0 \\
Llama-3-70B (base) & 30.2 & 30.2 & 31.8 & 33.9 & 35.1 & 37.4 \\
Llama-3-8B & 30.4 & 30.4 & 30.4 & 31.7 & 31.7 & 32.8 \\
MMed-Llama-3-8B & 36.8 & 36.8 & 38.6 & 40.2 & 40.2 & 42.6 \\
Mistral-7B & 37.1 & 37.1 & 37.1 & 36.3 & 36.3 & 36.3 \\
BioMistral-7B & 28.2 & 28.2 & 28.2 & 28.2 & 28.2 & 28.2 \\
}
\newcommand{\tableSmondalpha}{%
Qwen2.5-72B & 59.2 & 45.6 & 25.8 & 1.1 & 25.0 & 0.0 \\
HuatuoGPT-o1-72B & 65.0 & 44.1 & 27.3 & 3.1 & 25.0 & 0.0 \\
Med42-70B & 69.9 & 39.8 & 29.4 & 5.8 & 25.3 & 0.4 \\
OpenBioLLM-70B & 40.2 & 20.3 & 30.6 & 7.4 & 29.5 & 6.0 \\
MMed-Llama-3.1-70B & 41.3 & 21.7 & 35.2 & 13.6 & 26.6 & 2.1 \\
Llama-3.3-70B & 100.0 & 100.0 & 39.1 & 18.8 & 33.5 & 11.3 \\
Llama-3-70B (base) & 63.4 & 41.4 & 33.9 & 11.9 & 26.2 & 1.6 \\
Llama-3-8B & 61.8 & 39.0 & 31.7 & 9.0 & 25.8 & 1.0 \\
MMed-Llama-3-8B & 52.3 & 36.4 & 40.2 & 20.2 & 34.2 & 12.2 \\
Mistral-7B & 100.0 & 100.0 & 36.3 & 15.1 & 29.2 & 5.6 \\
BioMistral-7B & 40.5 & 20.7 & 28.2 & 4.3 & 25.5 & 0.6 \\
}
\newcommand{\tableSexact}{%
Qwen2.5-72B & 25.8 & 26.7 & +0.9 & NF4 & Yes \\
HuatuoGPT-o1-72B & 27.3 & 27.9 & +0.6 & NF4 & Yes \\
Med42-70B & 29.4 & 30.6 & +1.3 & NF4 & Yes \\
OpenBioLLM-70B & 30.6 & 31.5 & +1.0 & NF4 & Yes \\
MMed-Llama-3.1-70B & 35.2 & 35.2 & +0.0 & NF4 & Yes \\
Llama-3.3-70B & 39.1 & 40.0 & +0.9 & NF4 & Yes \\
Llama-3-70B (base) & 33.9 & 39.3 & +5.3 & NF4 & Yes \\
Llama-3-8B & 31.7 & 33.5 & +1.7 & NF4 & Yes \\
MMed-Llama-3-8B & 40.2 & 39.2 & -1.0 & NF4 & Yes \\
Mistral-7B & 36.3 & 36.3 & +0.0 & NF4 & Yes \\
BioMistral-7B & 28.2 & 28.2 & +0.0 & NF4 & Yes \\
}
\newcommand{\tableScalsize}{%
250 & 25.4 & 31.1 & 43.2 \\
500 & 25.0 & 27.9 & 40.3 \\
750 & 25.0 & 28.0 & 38.2 \\
1000 & 25.0 & 28.2 & 38.3 \\
1250 & 25.0 & 28.2 & 37.5 \\
}
\newcommand{\tableStaskwise}{%
RTN-8 & 100.0 & 1.1 & 100.0 & 1.4 \\
RTN-6 & 99.7 & 2.0 & 94.7 & 2.6 \\
RTN-4 & 60.0 & 1.9 & 56.5 & 3.7 \\
INT8 & 74.5 & 1.3 & 68.7 & 3.0 \\
NF4 & 94.6 & 2.3 & 91.5 & 3.8 \\
}
\newcommand{\tableSner}{%
\multicolumn{7}{l}{\textit{BC5CDR-Disease}}\\
Qwen2.5-72B & 0.723 & 0.722 & 0.733 & 0.669 & 0.711 & 0.956/0.566 \\
HuatuoGPT-o1-72B & 0.559 & 0.592 & 0.656 & 0.623 & 0.486 & 0.950/0.327 \\
Med42-70B & 0.863 & 0.856 & 0.847 & 0.693 & 0.823 & 0.885/0.769 \\
OpenBioLLM-70B & 0.707 & 0.705 & 0.713 & 0.684 & 0.572 & 0.918/0.415 \\
MMed-Llama-3.1-70B & 0.472 & 0.489 & 0.494 & 0.300 & 0.529 & 0.586/0.482 \\
Llama-3.3-70B & 0.866 & 0.877 & 0.883 & 0.843 & 0.860 & 0.947/0.788 \\
Llama-3-70B (base) & 0.825 & 0.819 & 0.816 & 0.816 & 0.856 & 0.951/0.778 \\
Llama-3-8B & 0.648 & 0.639 & 0.660 & 0.720 & 0.614 & 0.932/0.458 \\
MMed-Llama-3-8B & 0.453 & 0.462 & 0.452 & 0.433 & 0.418 & 0.397/0.443 \\
Mistral-7B & 0.577 & 0.594 & 0.599 & 0.601 & 0.620 & 0.770/0.519 \\
BioMistral-7B & 0.397 & 0.384 & 0.408 & 0.418 & 0.423 & 0.506/0.364 \\
\multicolumn{7}{l}{\textit{NCBI-Disease}}\\
Qwen2.5-72B & 0.745 & 0.743 & 0.747 & 0.745 & 0.740 & 0.957/0.604 \\
HuatuoGPT-o1-72B & 0.661 & 0.660 & 0.668 & 0.667 & 0.620 & 0.951/0.460 \\
Med42-70B & 0.722 & 0.718 & 0.705 & 0.642 & 0.608 & 0.893/0.460 \\
OpenBioLLM-70B & 0.577 & 0.576 & 0.578 & 0.501 & 0.551 & 0.946/0.389 \\
MMed-Llama-3.1-70B & 0.279 & 0.246 & 0.265 & 0.072 & 0.356 & 0.527/0.269 \\
Llama-3.3-70B & 0.717 & 0.716 & 0.723 & 0.756 & 0.690 & 0.951/0.541 \\
Llama-3-70B (base) & 0.678 & 0.675 & 0.670 & 0.685 & 0.676 & 0.940/0.528 \\
Llama-3-8B & 0.667 & 0.677 & 0.667 & 0.667 & 0.681 & 0.951/0.530 \\
MMed-Llama-3-8B & 0.409 & 0.412 & 0.371 & 0.390 & 0.510 & 0.655/0.417 \\
Mistral-7B & 0.476 & 0.477 & 0.497 & 0.504 & 0.462 & 0.815/0.323 \\
BioMistral-7B & 0.445 & 0.448 & 0.488 & 0.412 & 0.433 & 0.623/0.332 \\
}
\newcommand{\tableSprev}{%
\textbf{Maximum over models} & marginal & 9.5 & 8.0 & 6.6 & 5.7 & 6.3 & 6.6 \\
 & Mondrian & 2.7 & 2.8 & 3.1 & 3.5 & 3.8 & 3.9 \\
\textbf{Mean over models} & marginal & 3.7 & 3.5 & 3.4 & 3.3 & 3.1 & 3.1 \\
 & Mondrian & 1.8 & 1.8 & 1.9 & 2.0 & 2.0 & 2.1 \\
\midrule
Qwen2.5-72B & marginal & 2.6 & 3.2 & 3.7 & 4.2 & 4.7 & 4.9 \\
 & Mondrian & 1.8 & 2.1 & 2.4 & 2.7 & 3.1 & 3.2 \\
HuatuoGPT-o1-72B & marginal & 3.1 & 3.4 & 3.8 & 4.1 & 4.5 & 4.7 \\
 & Mondrian & 1.5 & 1.6 & 1.6 & 1.7 & 1.8 & 1.8 \\
Med42-70B & marginal & 4.0 & 3.8 & 3.7 & 3.5 & 3.4 & 3.3 \\
 & Mondrian & 2.7 & 2.7 & 2.7 & 2.8 & 2.8 & 2.9 \\
OpenBioLLM-70B & marginal & 9.5 & 8.0 & 6.6 & 5.2 & 3.8 & 3.1 \\
 & Mondrian & 2.4 & 2.4 & 2.4 & 2.4 & 2.4 & 2.4 \\
MMed-Llama-3.1-70B & marginal & 3.7 & 3.2 & 2.6 & 2.1 & 1.6 & 1.4 \\
 & Mondrian & 1.1 & 1.1 & 1.1 & 1.1 & 1.1 & 1.1 \\
Llama-3.3-70B & marginal & 2.0 & 2.4 & 2.8 & 3.2 & 3.6 & 3.8 \\
 & Mondrian & 1.0 & 1.1 & 1.3 & 1.5 & 1.7 & 1.8 \\
Llama-3-70B (base) & marginal & 3.9 & 4.5 & 5.1 & 5.7 & 6.3 & 6.6 \\
 & Mondrian & 2.5 & 2.8 & 3.1 & 3.5 & 3.8 & 3.9 \\
Llama-3-8B & marginal & 2.0 & 2.2 & 2.4 & 2.6 & 2.8 & 2.9 \\
 & Mondrian & 1.3 & 1.3 & 1.4 & 1.5 & 1.5 & 1.6 \\
MMed-Llama-3-8B & marginal & 4.6 & 3.7 & 2.7 & 1.8 & 1.0 & 0.5 \\
 & Mondrian & 1.9 & 1.8 & 1.7 & 1.6 & 1.6 & 1.5 \\
Mistral-7B & marginal & 3.2 & 2.6 & 2.0 & 1.4 & 0.7 & 0.3 \\
 & Mondrian & 1.0 & 0.9 & 0.7 & 0.6 & 0.5 & 0.4 \\
BioMistral-7B & marginal & 2.0 & 2.1 & 2.1 & 2.2 & 2.2 & 2.2 \\
 & Mondrian & 2.3 & 2.3 & 2.3 & 2.3 & 2.3 & 2.3 \\
}
   % \tableSdecdist, \tableSperconfig, \tableSaudit, \tableSmond, \tableSclassres,
                    % \tableSconfidence, \tableSexact, \tableScalsize, \tableSgrid, \tableSsplit, \tableSmondalpha,
                    % \tableSprev, \tableStaskwise, \tableSner

\title{\vspace{-3em}\textbf{Supplementary Information}\\[0.3em]
\large Do We Need the Largest LLMs for Every Clinical Decision?}
\author{}
\date{}

\begin{document}
\maketitle
\vspace{-3em}

\noindent This Supplementary Information provides the formal certificate, extended methods, complete
per-model audit tables, statistical and distributional robustness analyses, and the reproducibility
record for CERTA. It is organized as twelve sections: study design and task (S1); the formal
finite-sample guarantee with the certification algorithm (S2); the class-conditional derivation (S3);
the complete per-configuration and per-class certificate audit (S4); the two mechanisms by which a
decision class is escalated (S5); statistical robustness to confidence budget, calibration size,
threshold grid and split resampling (S6); distributional robustness to prevalence shift (S7);
deployment behaviour on escalated cases (S8); task-wise certification (S9); external human-gold
validation (S10); the deterministic-bound impossibility result (S11); and the reproducibility record of
models, software and commands (S12). The primary certificate and all main-text results use the
prespecified patient-disjoint split ($n_{\mathrm{cal}}=1{,}250$, $n_{\mathrm{test}}=1{,}250$,
$2{,}484$ patients); the robustness analyses below additionally repeat the procedure over twenty
independent patient-disjoint splits and re-run it with one decision per patient. All certificate and
robustness analyses use the recorded per-input reference decisions, quantized decisions and quantized
margins and require no additional model inference; the external named-entity results are taken from the
locked benchmark outputs. Confidence intervals for the clinical-verifier certificate are
patient-clustered bootstrap intervals (patients, not individual decisions, are resampled, respecting the
sixteen patients that contribute two decisions each); confidence intervals for the external
named-entity benchmarks (BC5CDR-Disease, NCBI-Disease), whose evaluation units are public documents
rather than patients, are document-level bootstrap intervals. All bootstraps use $1{,}000$ resamples.

\vspace{0.6em}
\noindent\fbox{\parbox{0.97\textwidth}{\small
\textbf{What CERTA does, and does not, certify.} CERTA controls \emph{compression-induced disagreement
risk}: the rate at which the quantized model's decision differs from the BF16 reference on the inputs
it is permitted to serve. It does \emph{not} control \emph{clinical-correctness risk}: whether the
reference decision is itself correct against a clinical ground truth. The external named-entity
benchmarks estimate a third, separate quantity, \emph{task-performance change} against human labels.
These three risks are kept distinct throughout: the certificate bounds the first, the benchmarks probe
the third, and the second is the object of the future clinician-adjudication study.}}

\vspace{0.6em}
\noindent\textbf{Clinical reading of the decision classes.} Operationally, $\YES$ accepts a candidate
clinical entity or answer and $\NO$ rejects it. A disagreement with the reference can correspond either
to a false accept or a false reject, whose clinical consequences may differ by workflow. CERTA
therefore does not impose an \emph{a priori} stricter risk on one class: the same target $\alpha$ is
controlled separately for the $\YES$ and $\NO$ classes, and whichever class cannot be certified at that
target is escalated. Empirically it is usually the reject ($\NO$) class that is escalated, because it is
the rarer decision and shows poorer compressed-versus-reference agreement (Supplementary
Tables~\ref{stab:mond} and~\ref{stab:classres}).

\tableofcontents
\vspace{1em}

% =====================================================================================
\section{Study design and task}
% =====================================================================================
Table~\ref{stab:notation} defines the quantities used throughout. The de-identified clinical narrative
corpus is the publicly available PMC-Patients dataset of patient summaries curated from open-access
PubMed Central case reports; it is sampled into the verifier accept/reject decision set that CERTA
certifies, this set is split patient-disjointly into calibration and held-out test, and the same model
checkpoints are separately evaluated on independent, human-annotated public named-entity benchmarks.

\begin{table}[h]\centering\small
\caption{\textbf{Interpretation of quantities.} Symbols and quantities used in the certificate and
throughout this Supplement.}
\label{stab:notation}
\setlength{\tabcolsep}{8pt}
\renewcommand{\arraystretch}{1.18}
\begin{tabular}{ll}
\toprule
Symbol & Meaning\\
\midrule
$x$ & a single decision input (a candidate clinical accept/reject decision)\\
$d^{\mathrm{ref}}(x)$ & bfloat16 (BF16) reference decision, $\YES$ (accept) or $\NO$ (reject)\\
$\hat d(x)$ & quantized-model decision on the same input\\
$g(x)$ & predicted decision class of the class-conditional certificate, $g(x)=\hat d(x)\in\{\YES,\NO\}$\\
$m_k(x)$ & decision margin $|z_\YES-z_\NO|$ at $k$-bit precision\\
$y(x)$ & disagreement indicator $\mathbf{1}[\hat d(x)\neq d^{\mathrm{ref}}(x)]$\\
$\alpha$ & target per-decision disagreement risk (risk budget)\\
$\delta$ & confidence budget; the certificate holds with probability $\ge 1-\delta$\\
$G$ & number of candidate thresholds in the Learn-then-Test grid\\
$\tau,\ \tau_g$ & selected margin threshold (marginal; per predicted class $g$)\\
$n_{\mathrm{cal}},\ n_{\mathrm{cal},g}$ & calibration inputs (total; in predicted class $g$)\\
$S_g(x;\tau_g)$ & class-$g$ service indicator, $\mathbf{1}[g(x)=g,\ m(x)\ge\tau_g]$\\
$\Risk(\tau)$ & marginal disagreement risk on the served region, $\Prob(\hat d\neq d^{\mathrm{ref}}\mid m\ge\tau)$\\
$\Risk_g(\tau_g)$ & class-$g$ disagreement risk, $\Prob(\hat d\neq d^{\mathrm{ref}}\mid g(x)=g,\ m\ge\tau_g)$\\
CP upper & one-sided Clopper--Pearson upper confidence bound on the risk\\
$\alpha_{\min,g}$ & zero-disagreement resolution floor $1-(\delta/2G)^{1/n_{\mathrm{cal},g}}$\\
coverage & fraction of decisions served by the quantized model\\
escalation & fraction routed to the BF16 reference ($=1-\text{coverage}$)\\
amortized size & coverage-weighted deployment cost relative to BF16 (\% BF16)\\
\bottomrule
\end{tabular}
\end{table}

The certified decision is the accept/reject step of an in-house clinical information-extraction
and verification pipeline that we have not yet published; only this accept/reject step is used or
evaluated here. Once the verifier decisions and quantized margins have been recorded, the certification
procedure itself does not depend on any pipeline-specific model internals; its distributional validity
depends only on how the inputs were sampled and on the exchangeability of calibration and test
decisions. In this pipeline, a generator model proposes candidate clinical entities and answers extracted from
de-identified patient narratives, and a verifier model adjudicates each candidate. The evaluated
decision space is binary, $\YES$ (accept) and $\NO$ (reject), and the model is constrained to a
single-token response rendered with each checkpoint's native chat template; for checkpoints whose
template does not accept a system role, the system instruction is merged into the user turn. Two
candidate types were evaluated in equal numbers: \emph{named-entity verification} (is a candidate
entity, with a proposed category, a valid and specific clinical entity of that category?) and
\emph{answer verification} (is a proposed answer supported by the supplied context and question?).
Reject cases were constructed by substituting a mismatched answer or a different clinical category, so
that the evaluation contains controlled reject cases (constructed from real clinical contexts) rather
than an accept-only distribution.

\paragraph{Decision distribution.} Because class imbalance in the verifier's own decisions is what
motivates the class-conditional policy, Supplementary Table~\ref{stab:decdist} reports, for every
model, the prevalence of $\YES$/$\NO$ decisions in the BF16 reference and in the NF4 quantized model,
overall and separately for the two task types. Several clinically adapted models are strongly
acceptance-skewed (for example the reject class of MMed-Llama-3.1-70B and OpenBioLLM-70B is a small
minority of their decisions), which is precisely the regime in which a single marginal threshold can
leave the rare reject decision poorly controlled.

\begin{table}[h]\centering\small
\caption{\textbf{Decision distribution.} Prevalence (\%) of accept ($\YES$) and reject ($\NO$)
decisions for the BF16 reference and the NF4 4-bit model, overall and split by task
(named-entity / answer verification), on the full evaluation set.}
\label{stab:decdist}
\setlength{\tabcolsep}{5pt}
\begin{tabular}{lrrrrcc}
\toprule
Model & \makecell{Ref\\YES} & \makecell{Ref\\NO} & \makecell{NF4\\YES} & \makecell{NF4\\NO}
& \makecell{Ref YES\\NER/QA} & \makecell{NF4 YES\\NER/QA}\\
\midrule
\tableSdecdist
\bottomrule
\end{tabular}
\end{table}

% =====================================================================================
\section{Formal CERTA guarantee}
% =====================================================================================
Algorithm~\ref{alg:certa} states the calibration procedure. To keep the candidate grid independent of
the data used to certify it, CERTA splits the calibration decisions into a \emph{threshold-selection}
subset, used only to construct the $G$-point candidate grid, and a disjoint \emph{certification} subset,
used only to evaluate the Clopper--Pearson bounds; coverage is then measured on the held-out test set.
Because the grid is fixed before the certification labels are read, the finite-sample Learn-then-Test
guarantee is exact: at each fixed candidate threshold the served certification decisions form a binomial
disagreement sample, a one-sided Clopper--Pearson upper bound at the Bonferroni-corrected level
$1-\delta/G$ controls its risk, and the smallest threshold whose bound clears the target is deployed (a
class with no admissible threshold is escalated in full). The class-conditional variant runs the same
procedure independently within each predicted decision class at $1-\delta/(2G)$. The certificate reported
in the main text uses the full $1{,}250$-decision calibration set for both the grid and the bounds, a
higher-power estimate; the split-sample certificate, which builds the grid on a $250$-decision
selection subset and evaluates the bounds on a disjoint $1{,}000$-decision certification subset (both
patient-disjoint from the $1{,}250$-decision test set), keeps the grid independent of the certification
labels and therefore carries the exact finite-sample guarantee of Proposition~\ref{prop:cert}. The
split-sample frontier reproduces the reported frontier to within a few percentage points and is
reported as an exactness sensitivity analysis in the statistical-robustness section.

\begin{algorithm}[t]
\caption{\textbf{CERTA certify-or-escalate calibration} (one decision class; run once per class for the
class-conditional variant).}
\label{alg:certa}
\small
\begin{flushleft}
\textbf{Input:} threshold-selection set $\mathcal{S}$ and certification set $\mathcal{C}$ (disjoint
calibration subsets), each a set of pairs $(m_i,y_i)$ with margin $m_i$ and disagreement flag
$y_i=\mathbf{1}[\hat d_i\neq d^{\mathrm{ref}}_i]$; target risk $\alpha$; confidence budget $\delta$; grid size $G$.\\
\textbf{Output:} threshold $\tau^\star$ and a serve/escalate rule.\\[0.35em]
$1.$\quad $\mathcal{T}\leftarrow$ the $G$-point candidate grid built from the selection set $\mathcal{S}$ (empirical margin quantiles)\hfill(fixed before $\mathcal{C}$ is read)\\
$2.$\quad $\delta'\leftarrow \delta/G$\hfill($\delta/(2G)$ per class for the class-conditional variant)\\
$3.$\quad \textbf{for} each $\tau\in\mathcal{T}$ \textbf{do}\\
$4.$\quad\quad $n_\tau\leftarrow|\{i\in\mathcal{C}:m_i\ge\tau\}|,\quad x_\tau\leftarrow|\{i\in\mathcal{C}:m_i\ge\tau,\ y_i=1\}|$\\
$5.$\quad\quad $U(\tau)\leftarrow \mathrm{BetaInv}(x_\tau+1,\ n_\tau-x_\tau,\ 1-\delta')$\hfill(one-sided Clopper--Pearson upper bound)\\
$6.$\quad $\tau^\star\leftarrow \min\{\tau\in\mathcal{T}:U(\tau)\le\alpha\}$; if the set is empty, $\tau^\star\leftarrow\infty$\hfill(escalate the whole class)\\
$7.$\quad \textbf{at deployment:} serve $x$ with the quantized model if $m(x)\ge\tau^\star$, otherwise escalate $x$ to the BF16 reference.
\end{flushleft}
\end{algorithm}

\begin{proposition}[Finite-sample decision-risk control]\label{prop:cert}
Fix $\alpha,\delta\in(0,1)$ and let the candidate grid $\mathcal{T}$ ($|\mathcal{T}|\le G$) be fixed
independently of the certification set, and suppose the certification and deployment decisions are
independent and identically distributed at the patient level. For each $\tau\in\mathcal{T}$ let $x_\tau$ and $n_\tau$ be the number of disagreements and
the number of served inputs in the certification set, and let
$\Risk(\tau)=\Prob(\hat d\neq d^{\mathrm{ref}}\mid m\ge\tau)$ denote the served-region disagreement risk.
The one-sided Clopper--Pearson upper bound $U(\tau)=\mathrm{BetaInv}(x_\tau+1,n_\tau-x_\tau,1-\delta/G)$ is
valid for $\Risk(\tau)$, and selecting the smallest $\tau^\star$ with $U(\tau^\star)\le\alpha$ gives
$\Prob\big(\Risk(\tau^\star)\le\alpha\big)\ge 1-\delta$ over the draw of the certification sample.
Running the procedure separately within each predicted class with budget $\delta/(2G)$ yields
$\Prob\big(\Risk_{\YES}\le\alpha \ \text{and}\ \Risk_{\NO}\le\alpha\big)\ge 1-\delta$.
\end{proposition}

\begin{proof}[Proof sketch]
Fix $\tau\in\mathcal{T}$. Since $\mathcal{T}$ is built from the selection set, which is disjoint from the
certification set, $\tau$ is independent of the certification labels; conditioning on the number of
served certification inputs $n_\tau$, the served-set disagreement count $x_\tau$ is
$\mathrm{Binomial}(n_\tau,\Risk(\tau))$ under the patient-level iid assumption, so the one-sided Clopper--Pearson quantity
$U(\tau)=\mathrm{BetaInv}(x_\tau+1,\,n_\tau-x_\tau,\,1-\delta/G)$ is a valid $1-\delta/G$ upper confidence
bound: $\Prob(\Risk(\tau)>U(\tau))\le\delta/G$. A Bonferroni union bound over the $|\mathcal{T}|\le G$
fixed candidates makes the probability that any candidate's bound understates its risk at most $\delta$;
on the complementary event the selected $\tau^\star$, which satisfies $U(\tau^\star)\le\alpha$ by
construction, obeys $\Risk(\tau^\star)\le U(\tau^\star)\le\alpha$, giving
$\Prob(\Risk(\tau^\star)\le\alpha)\ge 1-\delta$. For the class-conditional variant the same argument
applied within each predicted class at budget $\delta/(2G)$ gives each class a $1-\delta/2$ guarantee,
and a union bound over the two classes yields the joint statement.
\end{proof}

% =====================================================================================
\section{Class-conditional certification}
% =====================================================================================
The marginal certificate selects a single threshold $\tau$ by Learn-then-Test over a grid of
$G=50$ candidate thresholds, using a Bonferroni budget $\delta/G$ so that the finite-sample
disagreement bound $R_k(\tau)\le\alpha$ holds with confidence at least $1-\delta$. When two
class-conditional certificates are calibrated jointly (one per predicted decision class), controlling
each class independently at confidence $1-\delta$ would guarantee only $1-2\delta$ jointly. To retain
the stated $1-\delta$ over both classes, each class receives a Bonferroni budget of $\delta/(2G)$
across its own threshold grid, so that
\begin{equation}
\Prob\!\left(R_{\YES}\le\alpha \ \text{and}\ R_{\NO}\le\alpha\right)\ \ge\ 1-\delta ,
\end{equation}
by a union bound over the two class-conditional events. All class-conditional results in the main
text and in this Supplement use this $\delta/(2G)$ allocation. The correction is inexpensive, and it is a
\emph{distinct} comparison from the marginal-versus-class-conditional cost reported in the main text:
relative to the naive class-wise $\delta/G$ allocation, the joint $\delta/(2G)$ correction increases
each model's class-conditional frontier by at most $2.3$ percentage points of BF16 size, whereas the
main-text figure of $-0.8$ to $+7.9$ points (median $+2.3$) is the cost of the class-conditional
policy relative to the \emph{marginal} certificate. The flagship operating points (Qwen2.5-72B
$25.8\%$, HuatuoGPT-o1-72B $27.3\%$) are unchanged by the correction.

\paragraph{Sensitivity of the class-conditional frontier.} The class-conditional frontier inherits the
modest sensitivity of the marginal certificate to the confidence budget (Supplementary
Table~\ref{stab:conf}, Mondrian columns): the flagship operating points move by less than one
percentage point of BF16 size (Qwen2.5-72B $25.8\%$ to $26.2\%$, HuatuoGPT-o1-72B $27.3\%$ to
$28.2\%$) as confidence tightens from $90\%$ to $99\%$. Supplementary Table~\ref{stab:mondalpha}
reports the class-conditional frontier and its escalation fraction at $\alpha\in\{1,5,10\}\%$. Relaxing
the risk budget moves each model from a high-escalation regime (at $\alpha=1\%$ Qwen2.5-72B certifies
at $59.2\%$ of BF16 size and escalates $45.6\%$ of decisions) down toward the $4$-bit floor (at
$\alpha=10\%$ Qwen2.5-72B and HuatuoGPT-o1-72B reach $25.0\%$ with no escalation), with the deployed
$\alpha=5\%$ operating band ($25.8\%$ to $40.2\%$ across the eleven models) in between.

% =====================================================================================
\section{Complete certificate audit}
% =====================================================================================
Supplementary Table~\ref{stab:perconfig} gives the marginal certificate for every model and
configuration (round-to-nearest at $8$, $6$ and $4$ bits, $8$-bit LLM.int8() and NF4 $4$-bit): the
certified coverage and the empirical certified-set agreement at $\alpha=5\%$. Supplementary
Table~\ref{stab:audit} is the decision-risk audit for the deployed NF4 frontier points: for each
model it lists the number of certified calibration inputs, the number of those that disagree with the
reference, the selected margin threshold $\tau$, the one-sided Clopper--Pearson upper bound on the
certified-set risk, and the resulting certified coverage. Every Clopper--Pearson upper bound is at or
below the $5\%$ target, which is the finite-sample guarantee made auditable per operating point.
Supplementary Table~\ref{stab:mond} is the per-class decision-risk audit at every class-conditional
frontier point: for each decision class it lists the certified calibration count $n_g$, the number of
disagreements $x_g$, the selected threshold $\tau_g$, the certified coverage and the
Clopper--Pearson upper bound, so that each class certificate can be reconstructed independently.

\begin{table}[h]\centering\small
\caption{\textbf{Per-configuration marginal certificate.} Each cell is certified coverage (\%) with a
patient-clustered $95\%$ bootstrap confidence interval, at $\alpha=5\%$, $\delta=0.10$. The point
estimate is the certified test coverage and the brackets are patient-clustered $95\%$ bootstrap
intervals for that coverage; the formal risk certificate is the Clopper--Pearson upper bound
(Table~\ref{stab:audit}), not the bootstrap interval. INT8 fails on two BF16-native $70$B checkpoints,
which the certificate flags as $0\%$ coverage.}
\label{stab:perconfig}
\setlength{\tabcolsep}{4pt}
\resizebox{\textwidth}{!}{%
\begin{tabular}{lccccc}
\toprule
Model & RTN-8 & RTN-6 & RTN-4 & INT8 & NF4 \\
\midrule
\tableSperconfig
\bottomrule
\end{tabular}}
\end{table}

\begin{table}[h]\centering\small
\caption{\textbf{Decision-risk audit for the deployed NF4 frontier points ($\alpha=5\%$).}
$n_{\mathrm{cal}}$: certified calibration inputs ($m_k\ge\tau$); disag.: those disagreeing with the
BF16 reference; $\tau$: selected margin threshold; CP upper: one-sided Clopper--Pearson upper bound on
the certified-set disagreement risk (all $\le5\%$); coverage: certified test coverage. The
Clopper--Pearson upper bound is the formal certification quantity; the coverage and its $95\%$ bootstrap
CI are an empirical deployment estimate.}
\label{stab:audit}
\setlength{\tabcolsep}{7pt}
\begin{tabular}{lrrrrl}
\toprule
Model & $n_{\mathrm{cal}}$ & disag. & $\tau$ & \makecell{CP upper\\(\%)} & coverage (\%) [95\% CI] \\
\midrule
\tableSaudit
\bottomrule
\end{tabular}
\end{table}

\begin{table}[h]\centering\small
\caption{\textbf{Class-conditional (Mondrian) per-class decision-risk audit at the frontier
($\alpha=5\%$, $\delta/(2G)$).} For each decision class: $n$, certified calibration inputs;
$x$, disagreements among them; $\tau$, selected margin threshold ($\infty$ = class escalated in full);
Cov, certified test coverage; CP, one-sided Clopper--Pearson upper bound on the class risk (all
$\le5\%$). \emph{esc}: reject class escalated. Each class certificate is reconstructible from
$n$, $x$ and $\tau$.}
\label{stab:mond}
\setlength{\tabcolsep}{4pt}
\resizebox{\textwidth}{!}{%
\begin{tabular}{lrrrrr rrrrr}
\toprule
& \multicolumn{5}{c}{Accept ($\YES$)} & \multicolumn{5}{c}{Reject ($\NO$)}\\
\cmidrule(lr){2-6}\cmidrule(lr){7-11}
Model & $n_\YES$ & $x_\YES$ & $\tau_\YES$ & Cov (\%) & CP (\%)
      & $n_\NO$ & $x_\NO$ & $\tau_\NO$ & Cov (\%) & CP (\%)\\
\midrule
\tableSmond
\bottomrule
\end{tabular}}
\end{table}

% =====================================================================================
\section{Why classes are escalated}
% =====================================================================================
Whether a decision class can be certified at the target risk is governed by two distinct conditions,
and a class is escalated if it fails \emph{either}. The first is a resolution condition. With $x=0$
disagreements among $n_{\mathrm{cal},g}$ class-$g$ calibration inputs, the tightest certifiable class
risk is the zero-disagreement floor $\alpha_{\min,g}=1-(\delta/2G)^{1/n_{\mathrm{cal},g}}$; when
$\alpha_{\min,g}>\alpha$ the class has too few calibration inputs to reach the target under any
threshold and is escalated in full. Several clinically adapted models rarely return $\NO$, so their
reject class is a small minority of the calibration set (Supplementary Table~\ref{stab:classres}):
BioMistral-7B has only $4$ reject-class calibration inputs, giving a floor of $82.2\%$, and
OpenBioLLM-70B has $91$, giving a floor of $7.3\%$; both exceed the $5\%$ target, so their reject class
is escalated as \emph{resolution-limited}. The second is a risk condition. Even when the floor lies
below $\alpha$, the smallest Clopper--Pearson upper bound attainable over the threshold grid can still
exceed $\alpha$, because too many of the reject-class calibration decisions disagree with the reference
at every admissible threshold; the class is then escalated as \emph{risk-limited}.
MMed-Llama-3.1-70B is the illustrative case: its reject-class floor is $4.7\%$, below the target, yet
its best attainable Clopper--Pearson bound is $8.4\%$, so it is escalated despite having adequate
resolution. The zero-disagreement floor is therefore a \emph{necessary} resolution condition for
certification, not a \emph{sufficient} one. Supplementary Figure~\ref{sfig:classres} separates the two
mechanisms: every escalated model's reject class falls in either the resolution-limited region (floor
above $\alpha$) or the risk-limited region (floor below $\alpha$ but best attainable bound above it),
while the certified models occupy the low-floor, low-bound corner. This is the statistical reason for
the escalation behaviour in main-text Table~\ref{tab:frontier}, not an arbitrary choice.

\begin{center}\fbox{\parbox{0.9\textwidth}{\centering\small
The zero-disagreement floor is a \emph{necessary} resolution condition, not a \emph{sufficient}
certification condition.}}\end{center}

\begin{table}[h]\centering\small
\caption{\textbf{Class-specific calibration resolution and escalation mechanism (NF4 frontier).}
$n_\YES$, $n_\NO$: calibration inputs in each predicted class; $\alpha_{\min}$: zero-disagreement
resolution floor $1-(\delta/2G)^{1/n}$; min.\ CP$_\NO$: smallest Clopper--Pearson upper bound on the
reject-class risk attainable over the threshold grid; mechanism: whether the reject class is
\emph{certified}, \emph{resolution-limited} (floor $>\alpha$), or \emph{risk-limited} (floor
$\le\alpha$ but best attainable bound $>\alpha$). Target $\alpha=5\%$.}
\label{stab:classres}
\setlength{\tabcolsep}{6pt}
\begin{tabular}{lrrrrrl}
\toprule
Model & $n_\YES$ & $n_\NO$ & \makecell{$\alpha_{\min,\YES}$\\(\%)} & \makecell{$\alpha_{\min,\NO}$\\(\%)}
& \makecell{min.\ CP$_\NO$\\(\%)} & mechanism\\
\midrule
\tableSclassres
\bottomrule
\end{tabular}
\end{table}

% =====================================================================================
\section{Statistical robustness}
% =====================================================================================
\paragraph{Zero-disagreement floor.} With zero observed disagreements among $n_{\mathrm{cal}}$
calibration inputs, the tightest certifiable risk at confidence $1-\delta/G$ is
$1-(\delta/G)^{1/n_{\mathrm{cal}}}$, which at $\delta=0.10$, $G=50$, $n_{\mathrm{cal}}=1{,}250$ is
$\approx 0.5\%$; a $1\%$ target is therefore statistically reachable at this calibration size, and
targets below $\approx 0.5\%$ are not.

\paragraph{Confidence-budget sensitivity.} Supplementary Table~\ref{stab:conf} and Supplementary
Figure~\ref{sfig:robust}a report the marginal frontier at $\alpha=5\%$ for
$\delta\in\{0.10,0.05,0.01\}$ (confidence $90\%$, $95\%$, $99\%$), alongside the class-conditional
frontier at the same budgets. Across the evaluated confidence budgets the marginal frontier changed only
modestly: the range moves from $25.0$--$37.5\%$ at $\delta=0.10$ to $25.0$--$38.6\%$ at $\delta=0.01$,
and the median stays at $28.2\%$. The headline $25$--$40\%$ deployment envelope therefore does not depend on the choice of a
$90\%$ rather than a $95\%$ or $99\%$ confidence level.

\paragraph{Calibration-size sensitivity.} Supplementary Table~\ref{stab:calsize} and Supplementary
Figure~\ref{sfig:robust}b report the frontier as the calibration set is reduced from $1{,}250$ to
$250$ patient-disjoint decisions. The certificate degrades gracefully and monotonically: the least
compressible model moves from $37.5\%$ at $n_{\mathrm{cal}}=1{,}250$ to $43.2\%$ at
$n_{\mathrm{cal}}=250$, while the $4$-bit floor of $25\%$ is retained by the strongest models
throughout. More calibration data tightens the frontier, confirming that the $1\%$-resolution floor
is a statistical property of the calibration size rather than a model failure.

\paragraph{Threshold-grid sensitivity.} Learn-then-Test discretizes the margin axis into $G$ candidate
thresholds. Supplementary Table~\ref{stab:grid} and Supplementary Figure~\ref{sfig:gridsplit}a report
the marginal frontier as $G$ is swept over $\{25,50,100,200\}$. The frontier is essentially invariant:
the least-compressible model stays within $37.5$--$38.2\%$ of BF16 size, the median across models moves
only between $28.2\%$ and $29.4\%$, and the strongest models remain pinned at the $25\%$ $4$-bit floor
at every grid size. A finer grid tightens the per-candidate Bonferroni budget ($\delta/G$) but also
supplies more admissible thresholds; in the evaluated range these changes in threshold availability and
Bonferroni allocation produced little net change in the frontier, so $G=50$ is not a sensitive design
choice.

\paragraph{Repeated patient-disjoint splits.} A single calibration/test split could understate
variability. Supplementary Table~\ref{stab:split} and Supplementary Figure~\ref{sfig:gridsplit}b repeat
the entire certification over $20$ independent patient-disjoint splits, each re-shuffling the $2{,}484$
patient groups before the $50/50$ cut. The marginal frontier is stable across resamples: Qwen2.5-72B
sits exactly at the $25\%$ floor on all $20$ splits and HuatuoGPT-o1-72B within $0.8$ points of it,
while even the most variable model has an interquartile width of at most $2.5$ percentage points of
BF16 size. These results indicate that the observed $25$--$40\%$ deployment envelope is not driven by
the particular patient-disjoint split used for the primary analysis.

\paragraph{One decision per patient.} Sixteen of the $2{,}484$ patients contribute a second decision, so
the $2{,}500$-decision analysis carries a small amount of within-patient dependence. Re-running the
entire certification with one decision per patient (a single decision retained for each of the $2{,}484$
patients, split patient-disjointly) reproduces every frontier: the marginal frontier changes by at most
$1.8$ percentage points of BF16 size (MMed-Llama-3-8B) and the class-conditional frontier by at most
$1.1$ points (Llama-3-70B base), with the flagship operating points essentially unchanged (Qwen2.5-72B
$25.8\%$, HuatuoGPT-o1-72B $27.3\%$ to $27.6\%$). The primary analysis uses all $2{,}500$ decisions; the
one-decision-per-patient run is the formally patient-level exchangeable version of the procedure, and its
agreement with the full analysis (frontiers changing by at most $1.8$ percentage points) shows that the
retained within-patient dependence is immaterial to the reported results.

\paragraph{Risk-budget operating points.} Supplementary Table~\ref{stab:mondalpha} tabulates the
class-conditional deployment size and escalation fraction at $\alpha\in\{1,5,10\}\%$, giving the full
risk--resource trade-off for every model under the clinical default policy.

\paragraph{Split-conformal exactness check.} The reported certificate uses the full calibration set for
both the candidate grid and the Clopper--Pearson bounds, which maximizes statistical power but makes the
grid a function of the calibration data. The split-sample variant of Algorithm~\ref{alg:certa}
removes this dependence by building the grid on an independent $250$-decision selection subset and
evaluating the bounds on a disjoint $1{,}000$-decision certification subset, giving the exact
finite-sample guarantee of Proposition~\ref{prop:cert}. Re-running the class-conditional certificate in
this split-sample form reproduces the reported frontier closely (Supplementary Table~\ref{stab:exact}):
the deployment size changes by a median of $0.9$ and at most $5.3$ percentage points of BF16 size across
the eleven models, the flagship operating points move by less than one point (Qwen2.5-72B $25.8\%$ to
$26.7\%$, HuatuoGPT-o1-72B $27.3\%$ to $27.9\%$), and every model still selects the NF4 $4$-bit
configuration. The reported higher-power frontier is therefore not an artefact of the data-dependent grid.

\begin{table}[h]\centering\small
\caption{\textbf{Split-sample exactness.} Class-conditional deployment frontier (\% BF16, $\alpha=5\%$)
under the reported full-calibration certificate (grid and Clopper--Pearson bounds on all $1{,}250$
calibration decisions) and under the exact split-sample certificate (grid on a $250$-decision selection
subset, bounds on a disjoint $1{,}000$-decision certification subset). Difference is split minus full.
Every model selects the same NF4 $4$-bit deployable configuration and remains certified under the exact
split-sample procedure.}
\label{stab:exact}
\setlength{\tabcolsep}{7pt}
\begin{tabular}{lrrrcc}
\toprule
Model & \makecell{Full-cal.\\(\% BF16)} & \makecell{Split-sample\\(\% BF16)} & \makecell{Diff.\\(pp)} & \makecell{Frontier\\config} & \makecell{Exact split\\certified?}\\
\midrule
\tableSexact
\bottomrule
\end{tabular}
\end{table}

\begin{table}[h]\centering\small
\caption{\textbf{Confidence-budget sensitivity (marginal and class-conditional).} Frontier size
(\% BF16) at $\alpha=5\%$ for three confidence levels, for the marginal certificate and the
class-conditional (Mondrian) policy. Both change only modestly across $\delta$.}
\label{stab:conf}
\setlength{\tabcolsep}{7pt}
\begin{tabular}{lrrrrrr}
\toprule
& \multicolumn{3}{c}{Marginal} & \multicolumn{3}{c}{Class-conditional}\\
\cmidrule(lr){2-4}\cmidrule(lr){5-7}
Model & \makecell{$\delta{=}.10$\\(90\%)} & \makecell{$\delta{=}.05$\\(95\%)} & \makecell{$\delta{=}.01$\\(99\%)}
& \makecell{$\delta{=}.10$\\(90\%)} & \makecell{$\delta{=}.05$\\(95\%)} & \makecell{$\delta{=}.01$\\(99\%)}\\
\midrule
\tableSconfidence
\bottomrule
\end{tabular}
\end{table}

\begin{table}[h]\centering\small
\caption{\textbf{Calibration-size sensitivity.} Marginal frontier range (\% BF16) at $\alpha=5\%$
across the eleven models, as the patient-disjoint calibration set is reduced from $1{,}250$ to
$250$ decisions; the certificate degrades gracefully and the strongest models retain the $4$-bit floor.}
\label{stab:calsize}
\setlength{\tabcolsep}{10pt}
\begin{tabular}{rrrr}
\toprule
$n_{\mathrm{cal}}$ & \makecell{min\\($4\times$ floor)} & median & \makecell{max\\(least compressible)}\\
\midrule
\tableScalsize
\bottomrule
\end{tabular}
\end{table}

\begin{table}[h]\centering\small
\caption{\textbf{Threshold-grid sensitivity.} Marginal frontier size (\% BF16) at $\alpha=5\%$,
$\delta=0.10$ as the Learn-then-Test grid size $G$ is varied. The first three value columns are the
minimum, median and maximum frontier across the eleven models; the last three are three representative
models. The frontier is nearly invariant to $G$.}
\label{stab:grid}
\setlength{\tabcolsep}{7pt}
\begin{tabular}{rrrrrrr}
\toprule
$G$ & \makecell{min\\($4\times$ floor)} & median & \makecell{max\\(least compr.)} & \makecell{Qwen2.5\\72B} & \makecell{Huatuo\\o1-72B} & \makecell{Llama\\3.3-70B}\\
\midrule
\tableSgrid
\bottomrule
\end{tabular}
\end{table}

\begin{table}[h]\centering\small
\caption{\textbf{Split stability.} Marginal frontier size (\% BF16) at $\alpha=5\%$, $\delta=0.10$
across $20$ independent patient-disjoint calibration/test splits: per-model median, interquartile range
and full range. The frontier is stable across resamples; the strongest models stay at the $4$-bit floor
on every split.}
\label{stab:split}
\setlength{\tabcolsep}{9pt}
\begin{tabular}{lrrrr}
\toprule
Model & median & \makecell{IQR\\(25--75\%)} & min & max\\
\midrule
\tableSsplit
\bottomrule
\end{tabular}
\end{table}

\begin{table}[h]\centering\small
\caption{\textbf{Risk-budget operating points (class-conditional).} Class-conditional (Mondrian)
frontier size (\% BF16) and escalation fraction (\%) at three risk targets $\alpha\in\{1,5,10\}\%$
($\delta=0.10$, $\delta/(2G)$ allocation). Relaxing $\alpha$ lowers both the deployed size and the
escalation fraction toward the $4$-bit floor.}
\label{stab:mondalpha}
\setlength{\tabcolsep}{6pt}
\begin{tabular}{lrrrrrr}
\toprule
& \multicolumn{2}{c}{$\alpha=1\%$} & \multicolumn{2}{c}{$\alpha=5\%$} & \multicolumn{2}{c}{$\alpha=10\%$}\\
\cmidrule(lr){2-3}\cmidrule(lr){4-5}\cmidrule(lr){6-7}
Model & size & esc. & size & esc. & size & esc.\\
\midrule
\tableSmondalpha
\bottomrule
\end{tabular}
\end{table}

% =====================================================================================
\section{Distribution robustness: prevalence shift}
% =====================================================================================
Patient-disjoint calibration removes patient leakage but does not, by itself, guarantee robustness to
a shift between the calibration and deployment decision distributions. Because the marginal
certificate controls disagreement only on average over the calibration prevalence, a deployment whose
accept/reject mix differs from calibration can drive the marginal certified-set disagreement above the
target, whereas the class-conditional certificate controls each decision class separately and is
therefore insensitive to the tested class-prevalence shifts, provided within-class exchangeability
holds. Supplementary Figure~\ref{sfig:prevalence} and Supplementary Table~\ref{stab:prev} report the
certified-set disagreement as the deployment reject-class (NO) prevalence is swept from $5\%$ to $50\%$,
holding the calibrated thresholds fixed, for all eleven models. For OpenBioLLM-70B the marginal
certificate rises to $9.5\%$ at a balanced reject prevalence, well above the $5\%$ target, and for
Llama-3-70B and Qwen2.5-72B it drifts above $5\%$ at low reject prevalence. Taking the maximum over the
eleven models, the marginal certified-set disagreement exceeds the $5\%$ target at every deployment
prevalence ($5.7\%$ to $9.5\%$), whereas the class-conditional policy stays below the target for every
model at every prevalence (maximum over models $2.7\%$ to $3.9\%$). This supports adopting the
class-conditional policy as the default CERTA policy for the evaluated clinical verification setting.

\begin{table}[h]\centering\small
\caption{\textbf{Prevalence-shift stress test.} Certified-set disagreement (\%) versus deployment
reject-class (NO) prevalence, for the marginal and class-conditional (Mondrian) policies with fixed
calibrated thresholds. Values above the $\alpha=5\%$ target are the marginal certificate's
prevalence-induced violations; the class-conditional policy stays below target throughout.}
\label{stab:prev}
\setlength{\tabcolsep}{7pt}
\begin{tabular}{llrrrrrr}
\toprule
Model & Policy & \multicolumn{6}{c}{NO prevalence (\%)}\\
\cmidrule(lr){3-8}
 & & 50 & 40 & 30 & 20 & 10 & 5 \\
\midrule
\tableSprev
\bottomrule
\end{tabular}
\end{table}

% =====================================================================================
\section{Deployment behaviour: escalated cases}
% =====================================================================================
Supplementary Figure~\ref{sfig:escalated} contrasts the decisions the router serves with those it
escalates, at the NF4 frontier. Escalated inputs have an approximately $8.6$-fold smaller mean decision
margin ($0.93$ versus $7.95$ for certified inputs) and a much higher disagreement with the reference
(mean $45.3\%$ versus $2.9\%$). Escalation is therefore not arbitrary: the router preferentially refers low-margin decisions on which
disagreement with the reference is enriched, and serves the high-margin decisions on which the
compressed model is reliable. Whether those escalated disagreements are
clinically consequential is the object of the future clinician-adjudication study.

% =====================================================================================
\section{Task-wise certification}
% =====================================================================================
Supplementary Table~\ref{stab:taskwise} reports certified coverage and observed disagreement separately
for the two verifier tasks (named-entity and answer verification), averaged over models. The
certificate behaves consistently across both tasks at every precision (for example NF4 certifies
$94.6\%$ of named-entity and $91.5\%$ of answer decisions with $2.3\%$ and $3.8\%$ observed
disagreement, both within the $5\%$ target), showing consistent behaviour across the two evaluated
task types.

\begin{table}[h]\centering\small
\caption{\textbf{Task-wise certification.} Certified coverage and observed certified-set disagreement
(\%), averaged over the eleven models, for named-entity (NER) and answer (QA) verification at
$\alpha=5\%$, $\delta=0.10$.}
\label{stab:taskwise}
\setlength{\tabcolsep}{9pt}
\begin{tabular}{lrrrr}
\toprule
Configuration & \makecell{NER\\coverage} & \makecell{NER\\disag.} & \makecell{QA\\coverage} & \makecell{QA\\disag.}\\
\midrule
\tableStaskwise
\bottomrule
\end{tabular}
\end{table}

% =====================================================================================
\section{External human-gold validation}
% =====================================================================================
Supplementary Table~\ref{stab:ner} gives the full disease named-entity F1 against human annotations on
BC5CDR-Disease and NCBI-Disease, at BF16 and under round-to-nearest $8/6/4$-bit and the deployed NF4
$4$-bit configuration, together with precision and recall for the deployed NF4 configuration. These
benchmarks measure correctness against human labels, a different question from decision substitution;
the deployed NF4 configurations preserve F1 as well as or better than round-to-nearest $4$-bit (mean
absolute change $0.045$ on BC5CDR and $0.039$ on NCBI), and the domain model that collapses under
round-to-nearest $4$-bit (MMed-Llama-3.1-70B, $-0.172$ and $-0.207$) recovers under the deployed NF4
configuration ($+0.057$ and $+0.077$).

\begin{table}[h]\centering\small
\caption{\textbf{External disease named-entity results.} F1 against human annotations at BF16 and
under each configuration; the last column is precision/recall for the deployed NF4 $4$-bit
configuration.}
\label{stab:ner}
\setlength{\tabcolsep}{6pt}
\begin{tabular}{lrrrrrc}
\toprule
Model & \makecell{F1\\BF16} & \makecell{F1\\RTN-8} & \makecell{F1\\RTN-6} & \makecell{F1\\RTN-4}
& \makecell{F1\\NF4} & \makecell{NF4\\P/R}\\
\midrule
\tableSner
\bottomrule
\end{tabular}
\end{table}

% =====================================================================================
\section{Illustrative deterministic propagation bound}
% =====================================================================================
We evaluated whether a worst-case deterministic propagation of the quantization-induced logit
perturbation could bound the change in decision margin without any distributional assumption. Writing
$\Delta z(x)=z_k(x)-z_F(x)$ for the per-logit perturbation, a decision is guaranteed to be preserved
whenever the reference margin exceeds the worst-case margin perturbation,
$m_F(x) > \lVert \Delta z(x)\rVert_\infty$-type bounds propagated layerwise. The bound accumulates
multiplicatively through the residual and normalization structure as a product of per-block
operator-norm amplification factors. Here $c$ is a per-block worst-case amplification factor, that is,
an upper bound on the operator norm by which one transformer block can magnify a perturbation; for
attention and MLP blocks this factor generally exceeds one. Taking even the mild illustrative value
$c=1.5$ across $L\approx80$ blocks, the worst-case margin perturbation scales as $c^{L}$ and
exceeds $10^{14}$, whereas the observed decision margins are of order $10^{0}$--$10^{1}$ (median $5.1$;
Supplementary Fig.~\ref{sfig:escalated}a). Thus $c=1.5$ already yields a highly conservative illustrative
lower bound on the growth factor; the purpose of this calculation is to demonstrate the vacuity of naive
layerwise worst-case propagation at this scale, not to claim a tight operator-norm bound. The deterministic bound therefore exceeds every attainable margin by
more than a dozen orders of magnitude and certifies essentially no decisions. This vacuity is the
reason CERTA is built on the finite-sample, distribution-level Learn-then-Test argument of the main
text rather than on a deterministic per-input certificate; the deterministic route is reported here
only to document that it was considered and found uninformative at this scale.

% =====================================================================================
\section{Reproducibility}
% =====================================================================================
Supplementary Table~\ref{stab:inventory} lists the eleven evaluated checkpoints with their architecture
family, adaptation type and parameter count. Two independent quantization routes were evaluated so that
the frontier reflects a property of the certificate rather than of any single implementation:
round-to-nearest quantization at $8/6/4$ bits (symmetric per-output-channel scales over groups of $128$
elements, each group's scale chosen by a mean-squared-error search over clipping ratios
$\{0.7,0.8,0.9,1.0\}$, bfloat16 storage) and the bitsandbytes library (NF4 $4$-bit with double
quantization and bfloat16 compute, and $8$-bit LLM.int8()). The certificate itself is implemented in
Python with NumPy and SciPy: the Clopper--Pearson upper bound is the regularized incomplete
beta-function inverse (\texttt{scipy.special.betaincinv}), and confidence intervals are
patient-clustered bootstraps for the clinical-verifier certificate and document-level bootstraps for the
external benchmarks, from $1{,}000$ resamples. Given the recorded per-input reference decisions,
quantized decisions and quantized margins, every number in the main text and this Supplement is
regenerated by the following pipeline, which performs no model inference:

\begin{quote}\footnotesize\ttfamily
pip install -r requirements-lock.txt \ \# pinned dependency versions\\
python si\_compute.py \ \ \# per-config, audit, Mondrian, prevalence, task-wise, external NER\\
python si\_compute2.py \ \# patient-clustered bootstrap CIs, class resolution, margin ECDFs\\
python si\_compute3.py \ \# grid, split, class-conditional $\delta/\alpha$, escalation reason\\
python patient\_level\_check.py \# one-decision-per-patient robustness\\
python si\_tables.py \ \ \ \# emits si\_tables.tex\\
python build\_si\_figures.py \# emits figS1.pdf .. figS6.pdf
\end{quote}

\noindent The code release pins all dependency versions (an \texttt{environment.yml} and a
\texttt{requirements-lock.txt}) and is tagged with the exact commit used to produce the reported
results. The upstream inference that produced the recorded logits used the Hugging Face Transformers
library and PyTorch with no network access; the exact compute environment (GPU model and count, and
CUDA and driver versions) is specified by the authors in the code release, and NF4 and LLM.int8()
quantization require a CUDA accelerator.

\begin{table}[h]\centering\small
\caption{\textbf{Model and inference configuration.} The eleven evaluated checkpoints with
architecture/model family (Q, Qwen; L, Llama; M, Mistral), adaptation (G, general-purpose;
D, domain-adapted), parameter count, released checkpoint identifier,$^{\dagger}$ and measured peak VRAM
at the bfloat16 (BF16) reference and at the deployed NF4 $4$-bit configuration. Every model was
evaluated under the same quantizer set (round-to-nearest $8/6/4$-bit and the bitsandbytes NF4 $4$-bit
and LLM.int8() routes) with bfloat16 compute; the deployed frontier configuration is NF4 $4$-bit for all
eleven models.}
\label{stab:inventory}
\setlength{\tabcolsep}{5pt}
\resizebox{\textwidth}{!}{%
\begin{tabular}{llcrlcrr}
\toprule
Model & Family & Adapt. & \makecell{Params\\(B)} & HuggingFace checkpoint$^{\dagger}$ & \makecell{Compute\\dtype} & \makecell{BF16\\(GB)} & \makecell{NF4\\(GB)}\\
\midrule
Qwen2.5-72B & Qwen & G & 72 & \texttt{Qwen/Qwen2.5-72B-Instruct} & BF16 & 146.3 & 43.0\\
HuatuoGPT-o1-72B & Qwen & D & 72 & \texttt{FreedomIntelligence/HuatuoGPT-o1-72B} & BF16 & 146.1 & 42.7\\
Med42-70B & Llama & D & 70 & \texttt{m42-health/Llama3-Med42-70B} & BF16 & 141.9 & 41.2\\
OpenBioLLM-70B & Llama & D & 70 & \texttt{aaditya/Llama3-OpenBioLLM-70B} & BF16 & 141.7 & 41.0\\
Llama-3-70B (base) & Llama & G & 70 & \texttt{meta-llama/Meta-Llama-3-70B} & BF16 & 141.9 & 41.2\\
MMed-Llama-3.1-70B & Llama & D & 70 & \texttt{Henrychur/MMed-Llama3.1-70B} & BF16 & 141.9 & 41.2\\
Llama-3.3-70B & Llama & G & 70 & \texttt{meta-llama/Llama-3.3-70B-Instruct} & BF16 & 141.9 & 41.2\\
BioMistral-7B & Mistral & D & 7 & \texttt{BioMistral/BioMistral-7B} & BF16 & 14.8 & 4.6\\
Llama-3-8B & Llama & G & 8 & \texttt{meta-llama/Meta-Llama-3-8B} & BF16 & 16.7 & 6.4\\
Mistral-7B & Mistral & G & 7 & \texttt{mistralai/Mistral-7B-Instruct-v0.3} & BF16 & 14.9 & 4.7\\
MMed-Llama-3-8B & Llama & D & 8 & \texttt{Henrychur/MMed-Llama-3-8B} & BF16 & 16.7 & 8.7\\
\bottomrule
\end{tabular}}
\vspace{2pt}
\begin{minipage}{0.96\textwidth}\footnotesize
$^{\dagger}$Checkpoint identifiers give the released model families used; the authors confirm the exact
revisions in the code release. Peak VRAM is the measured allocation during inference; the same values
appear in main-text Table~\ref{tab:frontier}.
\end{minipage}
\end{table}

% =====================================================================================
% Supplementary Figures (collected; numbered S1--S7 in reading order)
% =====================================================================================
\clearpage
\section*{Supplementary Figures}

\begin{figure}[h]\centering
\includegraphics[width=0.86\textwidth]{figS1.pdf}
\caption{\textbf{Certified compression frontiers for all eleven models under the marginal
certificate.} Marginal deployment size (\% BF16) as a function of the risk budget $\alpha$; each curve
is one model, coloured by architecture family. The curves are smooth and monotonic and settle toward
the $4$-bit floor ($25\%$, dotted) as $\alpha$ relaxes. This is the full version of the representative
subset in main-text Fig.~\ref{fig:risk}a.}
\label{sfig:frontiers}
\end{figure}

\begin{figure}[h]\centering
\includegraphics[width=\textwidth]{figS2.pdf}
\caption{\textbf{Prevalence-shift stress test, all eleven models.} Certified-set disagreement versus
deployment reject-class (NO) prevalence, with calibrated thresholds held fixed, for the
\textbf{a}, marginal and \textbf{b}, class-conditional (Mondrian) policies. All eleven models are
shown (three highlighted, the other eight in grey). Under the marginal certificate several models
enter the shaded risk-budget violation zone above $\alpha=5\%$ (the maximum over models reaches
$9.5\%$); under the class-conditional policy every model stays below $\alpha$ at every prevalence
(maximum over models $3.9\%$; Supplementary Table~\ref{stab:prev}). The direction of the marginal
drift depends on which decision class carries more disagreement for each model.}
\label{sfig:prevalence}
\end{figure}

\begin{figure}[h]\centering
\includegraphics[width=\textwidth]{figS3.pdf}
\caption{\textbf{Escalated cases are enriched for low-margin, high-disagreement inputs
(NF4 $4$-bit, $\alpha=5\%$).}
\textbf{a}, Empirical cumulative distribution of the decision margin $m_k(x)$, pooled over the eleven
models, for escalated versus certified inputs; the distributions are strongly separated (median
margin $0.35$ escalated versus $5.12$ certified), not merely different in mean.
\textbf{b}, Observed disagreement with the BF16 reference for certified (filled) versus escalated
(open) inputs, per model: escalated inputs disagree with the reference far more often (mean $45.3\%$
versus $2.9\%$; dotted line, $5\%$). Escalation therefore selects the decisions on which the
compressed model is least reliable; whether those disagreements are clinically consequential is the
object of the future clinician-adjudication study.}
\label{sfig:escalated}
\end{figure}

\begin{figure}[h]\centering
\includegraphics[width=\textwidth]{figS4.pdf}
\caption{\textbf{Two mechanisms of reject-class escalation.}
\textbf{a}, Number of reject-class ($\NO$) calibration inputs per model (log scale), coloured by
escalation mechanism: certified (green), risk-limited (orange), resolution-limited (red).
Acceptance-skewed models have very few reject-class calibration inputs (BioMistral-7B, $4$;
OpenBioLLM-70B, $91$).
\textbf{b}, Reject-class resolution floor $\alpha_{\min,\NO}=1-(\delta/2G)^{1/n}$ (horizontal axis)
against the smallest attainable Clopper--Pearson upper bound on the reject-class risk (vertical axis),
both on log scales; dashed lines mark the $\alpha=5\%$ target. A class is certified only when both
coordinates fall below $\alpha$ (lower-left corner). Resolution-limited classes (BioMistral-7B,
OpenBioLLM-70B) fail the floor condition; the risk-limited class (MMed-Llama-3.1-70B) clears the floor
at $4.7\%$ but not the attainable bound at $8.4\%$. The floor is a necessary, not a sufficient,
certification condition.}
\label{sfig:classres}
\end{figure}

\begin{figure}[h]\centering
\includegraphics[width=\textwidth]{figS5.pdf}
\caption{\textbf{Statistical robustness.}
\textbf{a}, Marginal frontier size (\% BF16) at $\alpha=5\%$ for confidence budgets
$\delta\in\{0.10,0.05,0.01\}$ (confidence $90/95/99\%$); each line is one model. The frontier
changes only modestly with the confidence level.
\textbf{b}, Frontier range (minimum, median, maximum over models) at $\alpha=5\%$ as a function of the
calibration set size $n_{\mathrm{cal}}$; the certificate degrades gracefully as calibration data are
reduced and tightens as they increase, while the strongest models retain the $4$-bit floor.}
\label{sfig:robust}
\end{figure}

\begin{figure}[h]\centering
\includegraphics[width=\textwidth]{figS6.pdf}
\caption{\textbf{Design-choice robustness of the marginal frontier.}
\textbf{a}, Marginal frontier size (\% BF16) at $\alpha=5\%$ as the Learn-then-Test grid size $G$ is
varied over $\{25,50,100,200\}$: minimum ($4\times$ floor), median and maximum across the eleven
models. The envelope is nearly flat, so the frontier is insensitive to the grid resolution.
\textbf{b}, Distribution of the marginal frontier over $20$ independent patient-disjoint splits, one
boxplot per model (boxes coloured by architecture family, whiskers the full range; dotted line, the
$25\%$ $4$-bit floor). The strongest models sit at the floor on every split and no model's
interquartile width exceeds $2.5$ percentage points, confirming that the $25$--$40\%$ envelope is not
an artefact of the particular split used in the main text.}
\label{sfig:gridsplit}
\end{figure}

\end{document}
